\documentclass[UTF8,12pt]{ctexart}
\usepackage[a4paper,margin=24mm]{geometry}
\usepackage{amsmath,amssymb,bm,graphicx,adjustbox}
\newcommand{\fitmath}[1]{\begin{adjustbox}{max width=\linewidth}$\displaystyle #1$\end{adjustbox}}
\setlength{\parindent}{0pt}
\setlength{\parskip}{.7em}
\sloppy
\allowdisplaybreaks
\begin{document}
\section*{2017 年考研数学三 · 真题}
\subsection*{原 PDF 第 11 页}

\subsection*{2017年全国硕士研究生招生考试试题}

\subsection*{一、选择题（本题共8小题，每小题4分，共32分。在每小题给出的四个选项中，只有一项符合题目要求，把所选项前的字母填在题后的括号内。）}

（1）若函数 \(\displaystyle f(x)=\begin{cases}\dfrac{\displaystyle 1-\cos\sqrt{x}}{\displaystyle ax},\allowbreak &x>0,\allowbreak \\b,\allowbreak &x\le0\end{cases}\) 在 \(\displaystyle x=0\) 处连续，则（　）。


（A）\(\displaystyle ab=\dfrac{\displaystyle 1}{\displaystyle 2}\)。 （B）\(\displaystyle ab=-\dfrac{\displaystyle 1}{\displaystyle 2}\)。 （C）\(\displaystyle ab=0\)。 （D）\(\displaystyle ab=2\)。


（2）二元函数 \(\displaystyle z=xy(3-x-y)\) 的极值点是（　）。


（A）\(\displaystyle (0,\allowbreak 0)\)。 （B）\(\displaystyle (0,\allowbreak 3)\)。 （C）\(\displaystyle (3,\allowbreak 0)\)。 （D）\(\displaystyle (1,\allowbreak 1)\)。


（3）设函数 \(\displaystyle f(x)\) 可导，且 \(\displaystyle f(x)f^{\prime}(x)>0\)，则（　）。


（A）\(\displaystyle f(1)>f(-1)\)。 （B）\(\displaystyle f(1)<f(-1)\)。 （C）\(\displaystyle |f(1)|>|f(-1)|\)。 （D）\(\displaystyle |f(1)|<|f(-1)|\)。


（4）若级数 \(\displaystyle \sum_{n=2}^{\infty}[\sin\dfrac{\displaystyle 1}{\displaystyle n}-k\ln(1-\dfrac{\displaystyle 1}{\displaystyle n})]\) 收敛，则 \(\displaystyle k=\)（　）。


（A）\(\displaystyle 1\)。 （B）\(\displaystyle 2\)。 （C）\(\displaystyle -1\)。 （D）\(\displaystyle -2\)。


（5）设 \(\displaystyle \alpha\) 为 \(\displaystyle n\) 维单位列向量，\(\displaystyle E\) 为 \(\displaystyle n\) 阶单位矩阵，则（　）。


（A）\(\displaystyle E-\alpha\alpha^{\mathrm T}\) 不可逆。 （B）\(\displaystyle E+\alpha\alpha^{\mathrm T}\) 不可逆。 （C）\(\displaystyle E+2\alpha\alpha^{\mathrm T}\) 不可逆。 （D）\(\displaystyle E-2\alpha\alpha^{\mathrm T}\) 不可逆。


（6）已知矩阵 \(\displaystyle A=\begin{pmatrix}2&0&0\\0&2&1\\0&0&1\end{pmatrix}\)，\(\displaystyle B=\begin{pmatrix}2&1&0\\0&2&0\\0&0&1\end{pmatrix}\)，\(\displaystyle C=\begin{pmatrix}1&0&0\\0&2&0\\0&0&2\end{pmatrix}\)，则（　）。


（A）\(\displaystyle A\) 与 \(\displaystyle C\) 相似，\(\displaystyle B\) 与 \(\displaystyle C\) 相似。 （B）\(\displaystyle A\) 与 \(\displaystyle C\) 相似，\(\displaystyle B\) 与 \(\displaystyle C\) 不相似。 （C）\(\displaystyle A\) 与 \(\displaystyle C\) 不相似，\(\displaystyle B\) 与 \(\displaystyle C\) 相似。 （D）\(\displaystyle A\) 与 \(\displaystyle C\) 不相似，\(\displaystyle B\) 与 \(\displaystyle C\) 不相似。


（7）设 \(\displaystyle A,\allowbreak B,\allowbreak C\) 为三个随机事件，且 \(\displaystyle A\) 与 \(\displaystyle C\) 相互独立，\(\displaystyle B\) 与 \(\displaystyle C\) 相互独立，则 \(\displaystyle A\cup B\) 与 \(\displaystyle C\) 相互独立的充分必要条件是（　）。


（A）\(\displaystyle A\) 与 \(\displaystyle B\) 相互独立。 （B）\(\displaystyle A\) 与 \(\displaystyle B\) 互不相容。 （C）\(\displaystyle AB\) 与 \(\displaystyle C\) 相互独立。 （D）\(\displaystyle AB\) 与 \(\displaystyle C\) 互不相容。


（8）设 \(\displaystyle X_1,\allowbreak X_2,\allowbreak \cdots,\allowbreak X_n\;(n\ge2)\) 为来自总体 \(\displaystyle N(\mu,\allowbreak 1)\) 的简单随机样本，记 \(\displaystyle \bar X=\dfrac{\displaystyle 1}{\displaystyle n}\displaystyle \sum_{i=1}^nX_i\)，则下列结论中不正确的是（　）。


（A）\(\displaystyle \sum_{i=1}^n(X_i-\mu)^2\) 服从 \(\displaystyle \chi^2\) 分布。 （B）\(\displaystyle 2(X_n-X_1)^2\) 服从 \(\displaystyle \chi^2\) 分布。 （C）\(\displaystyle \sum_{i=1}^n(X_i-\bar X)^2\) 服从 \(\displaystyle \chi^2\) 分布。 （D）\(\displaystyle n(\bar X-\mu)^2\) 服从 \(\displaystyle \chi^2\) 分布。


\clearpage

\subsection*{原 PDF 第 12 页}

\subsection*{二、填空题（本题共6小题，每小题4分，共24分，把答案填在题中横线上。）}

（9）\(\displaystyle \int_{-\pi}^{\pi}(\sin^3x+\sqrt{\pi^2-x^2})\,\mathrm dx=\)\_\_\_\_\_\_。


（10）差分方程 \(\displaystyle y_{t+1}-2y_t=2^t\) 的通解为 \(\displaystyle y_t=\)\_\_\_\_\_\_。


（11）设生产某产品的平均成本为 \(\displaystyle \bar C(Q)=1+e^{-Q}\)，其中产量为 \(\displaystyle Q\)，则边际成本为\_\_\_\_\_\_。


（12）设函数 \(\displaystyle f(x,\allowbreak y)\) 具有一阶连续偏导数，且 \(\displaystyle \mathrm df(x,\allowbreak y)=ye^y\,\mathrm dx+x(1+y)e^y\,\mathrm dy\)，\(\displaystyle f(0,\allowbreak 0)=0\)，则 \(\displaystyle f(x,\allowbreak y)=\)\_\_\_\_\_\_。


（13）设矩阵 \(\displaystyle A=\begin{pmatrix}1&0&1\\1&1&2\\0&1&1\end{pmatrix}\)，\(\displaystyle \alpha_1,\allowbreak \alpha_2,\allowbreak \alpha_3\) 为线性无关的3维列向量组，则向量组 \(\displaystyle A\alpha_1,\allowbreak A\alpha_2,\allowbreak A\alpha_3\) 的秩为\_\_\_\_\_\_。


（14）设随机变量 \(\displaystyle X\) 的概率分布为 \(\displaystyle P\{X=-2\}=\dfrac{\displaystyle 1}{\displaystyle 2}\)，\(\displaystyle P\{X=1\}=a\)，\(\displaystyle P\{X=3\}=b\)，若 \(\displaystyle E(X)=0\)，则 \(\displaystyle D(X)=\)\_\_\_\_\_\_。


\subsection*{三、解答题（本题共9小题，共94分，解答应写出文字说明、证明过程或演算步骤。）}

（15）（本题满分10分）求 \(\displaystyle \lim_{x\to0^+}\dfrac{\displaystyle \int_0^x\sqrt{x-t}\,e^t\,\mathrm dt}{\displaystyle \sqrt{x^3}}\)。


（16）（本题满分10分）计算积分 \(\displaystyle \iint_D\dfrac{\displaystyle y^3}{\displaystyle (1+x^2+y^4)^2}\,\mathrm dx\mathrm dy\)，其中 \(\displaystyle D\) 是第一象限中以曲线 \(\displaystyle y=\sqrt{x}\) 与 \(\displaystyle x\) 轴为边界的无界区域。


\clearpage

\subsection*{原 PDF 第 13 页}

（17）（本题满分10分）求 \(\displaystyle \lim_{n\to\infty}\displaystyle \sum_{k=1}^n\dfrac{\displaystyle k}{\displaystyle n^2}\ln(1+\dfrac{\displaystyle k}{\displaystyle n})\)。


（18）（本题满分10分）已知方程 \(\displaystyle \dfrac{\displaystyle 1}{\displaystyle \ln(1+x)}-\dfrac{\displaystyle 1}{\displaystyle x}=k\) 在区间 \(\displaystyle (0,\allowbreak 1)\) 内有实根，确定常数 \(\displaystyle k\) 的取值范围。


（19）（本题满分10分）若 \(\displaystyle a_0=1,\allowbreak a_1=0,\allowbreak a_{n+1}=\dfrac{\displaystyle 1}{\displaystyle n+1}(na_n+a_{n-1})\;(n=1,\allowbreak 2,\allowbreak 3,\allowbreak \cdots)\)，\(\displaystyle S(x)\) 为幂级数 \(\displaystyle \sum_{n=0}^{\infty}a_nx^n\) 的和函数。（I）证明该幂级数的收敛半径不小于1。（II）证明 \(\displaystyle (1-x)S^{\prime}(x)-xS(x)=0\;(x\in(-1,\allowbreak 1))\)，并求 \(\displaystyle S(x)\) 的表达式。


（20）（本题满分11分）设3阶矩阵 \(\displaystyle A=(\alpha_1,\allowbreak \alpha_2,\allowbreak \alpha_3)\) 有3个不同的特征值，且 \(\displaystyle \alpha_3=\alpha_1+2\alpha_2\)。（I）证明 \(\displaystyle r(A)=2\)；（II）若 \(\displaystyle \beta=\alpha_1+\alpha_2+\alpha_3\)，求方程组 \(\displaystyle Ax=\beta\) 的通解。


\clearpage

\subsection*{原 PDF 第 14 页}

（21）（本题满分11分）设二次型 \(\displaystyle f(x_1,\allowbreak x_2,\allowbreak x_3)=2x_1^2-x_2^2+ax_3^2+2x_1x_2-8x_1x_3+2x_2x_3\) 在正交变换 \(\displaystyle x=Qy\) 下的标准形为 \(\displaystyle \lambda_1y_1^2+\lambda_2y_2^2\)，求 \(\displaystyle a\) 的值及一个正交矩阵 \(\displaystyle Q\)。


（22）（本题满分11分）设随机变量 \(\displaystyle X,\allowbreak Y\) 相互独立，且 \(\displaystyle X\) 的概率分布为 \(\displaystyle P\{X=0\}=P\{X=2\}=\dfrac{\displaystyle 1}{\displaystyle 2}\)，\(\displaystyle Y\) 的概率密度为 \(\displaystyle f(y)=\begin{cases}2y,\allowbreak &0<y<1,\allowbreak \\0,\allowbreak &\text{其他}.\end{cases}\)（I）求 \(\displaystyle P\{Y\le E(Y)\}\)；（II）求 \(\displaystyle Z=X+Y\) 的概率密度。


（23）（本题满分11分）某工程师为了解一台天平的精度，用该天平对一物体的质量做 \(\displaystyle n\) 次测量，该物体的质量 \(\displaystyle \mu\) 是已知的。设 \(\displaystyle n\) 次测量结果 \(\displaystyle X_1,\allowbreak X_2,\allowbreak \cdots,\allowbreak X_n\) 相互独立且均服从正态分布 \(\displaystyle N(\mu,\allowbreak \sigma^2)\)。该工程师记录的是 \(\displaystyle n\) 次测量的绝对误差 \(\displaystyle Z_i=|X_i-\mu|\;(i=1,\allowbreak 2,\allowbreak \cdots,\allowbreak n)\)，利用 \(\displaystyle Z_1,\allowbreak Z_2,\allowbreak \cdots,\allowbreak Z_n\) 估计 \(\displaystyle \sigma\)。（I）求 \(\displaystyle Z_1\) 的概率密度；（II）利用一阶矩求 \(\displaystyle \sigma\) 的矩估计量；（III）求 \(\displaystyle \sigma\) 的最大似然估计量。


\clearpage
\end{document}
